\chapter{Skin infections}
\medskip
\noindent\textit{Educational draft; seek professional advice for individual care.}

\section*{Definition and scope}
Skin infections is an infection-related topic requiring identification of the affected site, likely organism, severity, and transmission risk. Treatment depends on those details.

\section*{Contributors and risk}
Potential causes include bacteria, viruses, fungi, or parasites; exposure, travel, immune status, wounds, devices, and resistance patterns influence risk.

\section*{Presentation}
Record onset, fever, pain, rash, discharge, breathing or neurological symptoms, exposures, travel, vaccination, immune status, and contacts.

\section*{Assessment}
Use examination with targeted swabs, cultures, molecular tests, blood tests, or imaging when indicated. Not every infection needs antibiotics.

\section*{Management}
Care may include fluids, rest, isolation or notification advice, wound care, antimicrobials when indicated, and complication treatment.

\section*{When to seek urgent care}
Seek urgent care for sepsis features, breathing difficulty, confusion, severe dehydration, rapidly spreading redness, meningitis symptoms, or deterioration in an infant or immunocompromised person.

\section*{Prevention and follow-up}
Use vaccination, hand hygiene, safer sex, food and water safety, protective equipment, and appropriate contact precautions.

\section*{Expanded clinical context}
This chapter addresses skin infections, a topic that should be interpreted in the context of age, baseline health, medicines, comorbidities, goals, and access to care. It supports discussion with a qualified clinician and does not establish a diagnosis.

\section*{Assessment and decision points}
Review onset, duration, triggers, relevant exposures, physical findings, associated symptoms, functional impact, and immediate safety. Record the main concern, time course, risk factors, red flags, and the person’s questions. Use targeted investigations when their result is likely to change management.

\section*{Management and follow-up}
Care may include education, practical support, treatment of contributing conditions, medication or a procedure when indicated, and a shared follow-up plan. Explain expected improvement, possible adverse effects, and when reassessment is needed. Avoid starting, stopping, or sharing prescription medicines without professional advice.

\section*{Safety-netting}
Seek urgent help for severe or rapidly worsening symptoms, collapse, breathing difficulty, severe pain, confusion, dehydration, uncontrolled bleeding, or any immediate danger.

\section*{Practical prevention}
Use plain-language information, supportive routines, risk reduction, and professional review when symptoms persist, recur, or impair daily life. Keep written instructions and confirm how to access help outside routine appointments.

\section*{Limitations}
This educational expansion should be checked against current local guidance and authoritative topic-specific sources.
\section*{Sources and limitations}
This concise educational overview should be checked against current local clinical guidance and adapted to the individual.
